\chapter{Entropy, KL Divergence, and Fisher Information}
\chaptermark{Entropy and Fisher Information}
\label{chap:information}

\begingroup
\clubpenalty=10000
\widowpenalty=10000
% Shared scientific charts for the notes and Lecture 7 slides.
\providecommand{\InfoDataPath}{figures/data}
\providecommand{\InfoPlotWidth}{11.8cm}
\providecommand{\InfoPlotHeight}{4.4cm}
\definecolor{InfoBlue}{HTML}{1F4E79}
\definecolor{InfoOrange}{HTML}{C46A13}
\definecolor{InfoTeal}{HTML}{007D80}
\definecolor{InfoMuted}{HTML}{526477}
\pgfplotsset{info axis/.style={
  scale only axis,width=\InfoPlotWidth,height=\InfoPlotHeight,
  axis lines=left,enlargelimits=false,
  grid=major,grid style={black!9},
  tick label style={font=\footnotesize},label style={font=\footnotesize},
  /pgf/number format/.cd,fixed,precision=2,/pgfplots/.cd,
  legend style={draw=none,fill=none,font=\footnotesize,
    at={(0.5,1.04)},anchor=south,legend columns=-1,column sep=0.7em},
  every axis plot/.append style={no marks,line width=1.15pt},
}}

\newcommand{\InfoEntropyPlot}{%
\begin{tikzpicture}\begin{axis}[info axis,
  xmin=0,xmax=1,ymin=0,ymax=0.74,
  xlabel={Success probability $p$},ylabel={Entropy $H(P_p)$, nats},
  xtick={0,0.2,0.5,0.8,1},ytick={0,0.2,0.4,0.693147},
  yticklabels={$0$,$0.2$,$0.4$,$\log 2$}]
  \addplot[InfoBlue] table[x=p,y=entropy] {\InfoDataPath/ch7-entropy.dat};
  \addplot[InfoOrange,only marks,mark=*,mark size=2pt]
    coordinates {(0.5,0.69314718056) (0.8,0.500402423538)};
\end{axis}\end{tikzpicture}%
}

\newcommand{\InfoCrossEntropyPlot}{%
\begin{tikzpicture}\begin{axis}[info axis,
  xmin=0.02,xmax=0.98,ymin=0,ymax=3.3,
  xlabel={Predicted success probability $q$},ylabel={Expected loss, nats},
  xtick={0.02,0.2,0.5,0.8,0.98},ytick={0,0.5,1,2,3},
  scaled ticks=false]
  \addplot[InfoBlue] table[x=q,y=ce] {\InfoDataPath/ch7-cross-entropy.dat};
  \addlegendentry{$\operatorname{CE}(P,Q)$}
  \addplot[InfoOrange,densely dotted] table[x=q,y=entropy] {\InfoDataPath/ch7-cross-entropy.dat};
  \addlegendentry{$H(P)$}
  \addplot[InfoTeal,dashed] table[x=q,y=kl] {\InfoDataPath/ch7-cross-entropy.dat};
  \addlegendentry{$D_{\mathrm{KL}}(P\Vert Q)$}
  \addplot[InfoMuted,densely dotted,forget plot]
    coordinates {(0.8,0) (0.8,3.3)};
\end{axis}\end{tikzpicture}%
}

\newcommand{\InfoLikelihoodPlot}{%
\begin{tikzpicture}\begin{axis}[info axis,
  xmin=-2,xmax=2,ymin=0,ymax=1.04,
  xlabel={Candidate mean minus fitted mean $\mu-\bar x$},
  ylabel={Relative likelihood $L(\mu)/L(\bar x)$},
  xtick={-2,-1,0,1,2},ytick={0,0.25,0.5,0.75,1}]
  \addplot[InfoOrange] table[x=offset,y=n5] {\InfoDataPath/ch7-likelihood.dat};
  \addlegendentry{$n=5$}
  \addplot[InfoBlue,dashed] table[x=offset,y=n20] {\InfoDataPath/ch7-likelihood.dat};
  \addlegendentry{$n=20$}
  \addplot[InfoTeal] table[x=offset,y=n80] {\InfoDataPath/ch7-likelihood.dat};
  \addlegendentry{$n=80$}
\end{axis}\end{tikzpicture}%
}

\newcommand{\InfoFisherPlot}{%
\begin{tikzpicture}\begin{axis}[info axis,
  xmin=0.02,xmax=0.98,ymin=0,ymax=55,
  xlabel={Bernoulli parameter $p$},ylabel={Fisher information $I_1(p)$},
  xtick={0.02,0.2,0.5,0.8,0.98},ytick={4,10,20,40,50},
  scaled ticks=false]
  \addplot[InfoBlue] table[x=p,y=information] {\InfoDataPath/ch7-fisher.dat};
  \addplot[InfoOrange,only marks,mark=*,mark size=2pt]
    coordinates {(0.2,6.25) (0.5,4) (0.8,6.25)};
\end{axis}\end{tikzpicture}%
}

\newcommand{\InfoVariancePlot}{%
\begin{tikzpicture}\begin{loglogaxis}[info axis,
  xmin=1,xmax=200,ymin=0.015,ymax=6,
  xlabel={Independent observations $n$},ylabel={Estimator variance},
  xtick={1,5,20,100,200},xticklabels={$1$,$5$,$20$,$100$,$200$},
  ytick={0.02,0.1,1,4},yticklabels={$0.02$,$0.1$,$1$,$4$}]
  \addplot[InfoTeal] table[x=n,y=bound] {\InfoDataPath/ch7-variance.dat};
  \addlegendentry{Bound $4/n$}
  \addplot[InfoBlue,only marks,mark=*,mark size=1.8pt]
    table[x=n,y=meanvariance] {\InfoDataPath/ch7-variance.dat};
  \addlegendentry{Simulated $\bar X$}
  \addplot[InfoOrange,dashed] table[x=n,y=firsttheory] {\InfoDataPath/ch7-variance.dat};
  \addlegendentry{Using only $X_1$}
\end{loglogaxis}\end{tikzpicture}%
}

\newcommand{\InfoLocalKLPlot}{%
\begin{tikzpicture}\begin{axis}[info axis,
  xmin=0.65,xmax=0.95,ymin=0,ymax=0.15,
  xlabel={Nearby Bernoulli parameter $q$},ylabel={KL divergence, nats},
  xtick={0.65,0.7,0.8,0.9,0.95},ytick={0,0.05,0.1,0.15},
  scaled ticks=false]
  \addplot[InfoBlue] table[x=q,y=kl] {\InfoDataPath/ch7-local-kl.dat};
  \addlegendentry{Exact $D_{\mathrm{KL}}(P_{0.8}\Vert P_q)$}
  \addplot[InfoOrange,dashed] table[x=q,y=quadratic] {\InfoDataPath/ch7-local-kl.dat};
  \addlegendentry{$\tfrac12 I_1(0.8)(q-0.8)^2$}
\end{axis}\end{tikzpicture}%
}


% Mathematical background checked against Goodfellow et al., Deep Learning,
% Chapter 3, and Stanford STATS 200 (2016), Lectures 14--15.
% Examples, numerical tables, and figures are generated in this repository.

\section{Surprise and entropy}
\label{sec:entropy}

Suppose a population has two labels, with probabilities
\(\Prob(Y=1)=0.8\) and \(\Prob(Y=0)=0.2\).
Before seeing the next label, we expect \(1\). Seeing \(0\) is less
likely and tells us more than another \(1\). We want a numerical measure
of that difference.

Multiplying probabilities should add information. For example, two
independent outcomes of probability \(1/2\) have joint probability
\(1/4\). The function \(-\log p\) has exactly this property:
\[
  -\log(p_1p_2)=-\log p_1-\log p_2.
\]
It is zero for an outcome certain to occur and grows as the outcome
becomes less likely.

\begin{definition}[Surprisal]
If an observed outcome has probability \(p>0\), its surprisal is
\(-\log p\). Throughout this chapter \(\log\) is the natural logarithm,
so information is measured in \emph{nats}.
\end{definition}

For our binary example, the surprisals are
\[
  -\log0.8\approx0.2231,\qquad -\log0.2\approx1.6094.
\]
The rarer label has greater surprisal. To describe the uncertainty
before the label is observed, average these values using the actual
label probabilities.

\begin{definition}[Entropy of a finite distribution]
Let \(P\) assign probabilities \(p_1,\ldots,p_K\) to \(K\) labels.
Its entropy is the expected surprisal:
\begin{equation}
  H(P)=\Expect_{Y\sim P}[-\log p_Y]
      =-\sum_{k=1}^{K}p_k\log p_k.
  \label{eq:entropy}
\end{equation}
A term with \(p_k=0\) is defined to be zero, using
\(\lim_{p\downarrow0}p\log p=0\).
\end{definition}

Thus \(H(P)=-0.8\log0.8-0.2\log0.2\approx0.5004\) nats.
This is an average over repeated draws, not the surprisal of every draw.
If one label has probability one, the entropy is zero: there is no
uncertainty about the next label.

Write \(P_p\) for the Bernoulli distribution with success probability \(p\).
Then
\[
  H(P_p)=-p\log p-(1-p)\log(1-p).
\]
For \(0<p<1\), differentiating gives
\[
  \frac{d}{dp}H(P_p)=\log\frac{1-p}{p},\qquad
  \frac{d^2}{dp^2}H(P_p)=-\frac1{p(1-p)}<0.
\]
The first derivative vanishes at \(p=1/2\), so the largest entropy is
\(\log2\). A balanced binary outcome is harder to predict than an outcome
that is almost always the same.

\begin{figure}[H]
  \centering\InfoEntropyPlot
  \caption{Bernoulli entropy. The maximum occurs at \(p=1/2\).
  The second marked point is \(p=0.8\), where \(H(P_p)\approx0.5004\).
  Both deterministic endpoints have entropy zero.}
  \label{fig:binary-entropy}
\end{figure}

Changing the logarithm base changes only the unit. Dividing entropy in
nats by \(\log2\) gives entropy in bits. We keep nats so that the formulas
use the same logarithm as the likelihood calculations in Chapter~\ref{chap:estimation}.

\section{Cross-entropy and KL divergence}
\label{sec:cross-entropy-kl}

Suppose the actual label distribution is \(P=(p_1,\ldots,p_K)\), but a
predictor reports \(Q=(q_1,\ldots,q_K)\) for the same labels.
When label \(k\) occurs, its log loss is \(-\log q_k\).
We average using \(P\), because \(P\) describes which labels actually occur.

\begin{definition}[Cross-entropy]
The cross-entropy of \(P\) relative to predictions \(Q\) is
\begin{equation}
  \operatorname{CE}(P,Q)=-\sum_{k=1}^Kp_k\log q_k.
  \label{eq:cross-entropy}
\end{equation}
Terms with \(p_k=0\) contribute zero. If \(p_k>0\) but \(q_k=0\),
the cross-entropy is infinite: a possible outcome was declared impossible.
\end{definition}

For the population with success probability \(0.8\), a predictor assigning
success probability \(q\) has expected loss
\[
  \operatorname{CE}(P_{0.8},P_q)
  =-0.8\log q-0.2\log(1-q).
\]
\Needspace{8\baselineskip}
The following values compare predictions for the same population:
\begin{center}
  \begin{tabular}{@{}rrr@{}}
    \toprule
    Predicted \(q\) & Cross-entropy & Excess over \(H(P_{0.8})\)\\
    \midrule
    \(0.50\) & \(0.6931\) & \(0.1927\)\\
    \(0.60\) & \(0.5919\) & \(0.0915\)\\
    \(0.80\) & \(0.5004\) & \(0\)\\
    \(0.99\) & \(0.9291\) & \(0.4287\)\\
    \bottomrule
  \end{tabular}
\end{center}
Predicting \(0.99\) gives a small loss on the frequent successes but a
large loss on the failures. On average it is worse than predicting \(0.5\).
The best value is \(q=0.8\); even there, the expected loss is not zero
because individual labels remain random.

\begin{definition}[Kullback--Leibler divergence]
For distributions on the same finite set,
\begin{equation}
  D_{\mathrm{KL}}(P\Vert Q)
  =\sum_{k:p_k>0}p_k\log\frac{p_k}{q_k}.
  \label{eq:kl}
\end{equation}
The value is infinite if \(q_k=0\) for any \(k\) with \(p_k>0\).
Subtracting the entropy formula from the cross-entropy formula gives
\begin{equation}
  \operatorname{CE}(P,Q)=H(P)+D_{\mathrm{KL}}(P\Vert Q).
  \label{eq:ce-entropy-kl}
\end{equation}
\end{definition}

The direction matters: the expectation in \(D_{\mathrm{KL}}(P\Vert Q)\)
is over outcomes drawn from \(P\). Here KL measures the extra expected
log loss caused by using \(Q\) in place of the actual distribution.

\begin{result}[Nonnegativity]
\(D_{\mathrm{KL}}(P\Vert Q)\ge0\), with equality exactly when \(P=Q\).
Consequently, \(\operatorname{CE}(P,Q)\ge H(P)\).
\end{result}

\begin{proof}
If the divergence is infinite, the inequality holds. For \(u>0\), the
function \(u-1-\log u\) has derivative \(1-1/u\), so its minimum is zero
at \(u=1\). Thus \(-\log u\ge1-u\). Apply this to
\(u=q_k/p_k\) for every \(p_k>0\):
\[
  D_{\mathrm{KL}}(P\Vert Q)
  \ge\sum_{k:p_k>0}(p_k-q_k)
  =1-\sum_{k:p_k>0}q_k\ge0.
\]
The scalar inequality is an equality only at \(u=1\).
Equality therefore requires \(q_k=p_k\) on the support of \(P\);
normalization then leaves no mass for any other label.
\end{proof}

KL is not a distance. For example,
\[
  D_{\mathrm{KL}}(P_{0.8}\Vert P_{0.5})\approx0.1927,\qquad
  D_{\mathrm{KL}}(P_{0.5}\Vert P_{0.8})\approx0.2231.
\]
It is not symmetric and does not generally satisfy the triangle inequality.
When writing a KL penalty, we must specify its direction.

\begin{figure}[H]
  \centering\InfoCrossEntropyPlot
  \caption{The population probability stays \(0.8\), while the prediction
  \(q\) changes. The vertical gap between cross-entropy and the constant
  entropy is KL. Both cross-entropy and KL are minimized at \(q=0.8\).}
  \label{fig:cross-entropy-kl}
\end{figure}

The same result identifies the maximum categorical entropy. If \(U\)
is uniform on \(K\) labels, then
\[
  D_{\mathrm{KL}}(P\Vert U)=\log K-H(P)\ge0.
\]
Hence \(0\le H(P)\le\log K\), with the upper bound attained only by the
uniform distribution.

\section{Connections to learning}
\label{sec:information-learning}

\subsection{Likelihood and empirical cross-entropy}

Chapter~\ref{chap:estimation} derived log loss from likelihood.
We can now give that loss an interpretation. Suppose \(y_1,\ldots,y_n\)
are independent observations taking values in \(K\) labels, and a model
assigns probability \(q_\theta(k)\) to label \(k\). Let \(n_k\) be the
number of occurrences of label \(k\), and set \(\widehat p_k=n_k/n\).
Then
\begin{equation}
  -\frac1n\log L(\theta)
  =-\frac1n\sum_{i=1}^n\log q_\theta(y_i)
  =-\sum_{k=1}^K\widehat p_k\log q_\theta(k)
  =\operatorname{CE}(\widehat P,Q_\theta).
  \label{eq:mle-empirical-ce}
\end{equation}
The empirical distribution \(\widehat P\) is the one introduced in
Section~\ref{sec:bootstrap}. Its entropy is fixed once the data are fixed.
Therefore maximum likelihood minimizes
\(D_{\mathrm{KL}}(\widehat P\Vert Q_\theta)\) in this finite-label setting.
This does not claim that the fitted model equals the population distribution:
the empirical frequencies are noisy, and the model may impose restrictions.

For a classifier, predictions depend on input \(x\). Write \(q_\theta(k\mid x)\)
for its predicted probability of label \(k\).
For an observed pair \((x_i,y_i)\), define the \emph{one-hot} target by
\(t_{ik}=1\) if \(k=y_i\), and \(t_{ik}=0\) otherwise. Its cross-entropy is
\[
  -\sum_{k=1}^K t_{ik}\log q_\theta(k\mid x_i)
  =-\log q_\theta(y_i\mid x_i).
\]
Averaging over examples gives the multiclass negative log-likelihood.
For two labels, this is the binary log loss already derived in
Section~\ref{sec:likelihood-loss}.

The population expected loss also separates into entropy and KL for each
input, using the true conditional label distribution \(P(\,\cdot\mid x)\).
An individual one-hot target does not mean that this true conditional
distribution is deterministic.

\subsection{Sequences and soft targets}

A language model assigns a conditional distribution to the next token.
For a sequence \(w_1,\ldots,w_m\), the chain rule from
Chapter~\ref{chap:probability-models} gives
\[
  Q_\theta(w_1,\ldots,w_m)
  =\prod_{t=1}^m q_\theta(w_t\mid w_1,\ldots,w_{t-1}).
\]
At \(t=1\), the condition is empty. The average negative log-likelihood is
\[
  -\frac1m\sum_{t=1}^m
    \log q_\theta(w_t\mid w_1,\ldots,w_{t-1}).
\]
For three observed tokens given conditional probabilities
\(0.5,0.25,0.1\), the sequence probability is \(0.0125\), and the
average loss is \(-\log(0.0125)/3\approx1.4607\) nats per token.
The product follows from conditioning; the tokens need not be independent.

Targets can also be probability vectors. In \emph{distillation}, a
teacher model supplies probabilities \(P\), and a student supplies \(Q\)
for the same input. Minimizing \(\operatorname{CE}(P,Q)\) matches all
teacher probabilities, not just the most likely label. For example,
\[
  P=(0.7,0.2,0.1),\qquad Q=(0.6,0.25,0.15)
\]
give \(H(P)\approx0.8018\), \(\operatorname{CE}(P,Q)\approx0.8245\),
and \(D_{\mathrm{KL}}(P\Vert Q)\approx0.0227\).
The teacher entropy is constant with respect to the student parameters,
so minimizing this cross-entropy is equivalent to minimizing this KL.

\subsection{KL penalties and approximate posteriors}

A KL penalty can discourage a model from moving far from a reference
distribution. For example, consider
\[
  \mathcal L_{\mathrm{train}}(\theta)
  +\beta D_{\mathrm{KL}}(Q_\theta\Vert Q_{\mathrm{ref}}),
  \qquad \beta\ge0.
\]
Here \(\mathcal L_{\mathrm{train}}\) is the training loss,
\(Q_\theta\) is the model's output distribution on a specified input,
and \(Q_{\mathrm{ref}}\) is a fixed reference on the same labels.
The coefficient \(\beta\) sets the penalty strength. For many inputs,
one can average the penalty over inputs. The direction written here
penalizes probabilities used by \(Q_\theta\) that are small under the
reference. This is an output-distribution penalty; the MAP penalty in
Section~\ref{sec:map-estimation} instead came from a prior on parameters.

KL also gives a way to approximate a posterior. Consider a finite unknown
state \(Z\), observed data \(x\), and a joint model \(p(x,z)\).
Suppose the posterior \(P(z\mid x)=p(x,z)/p(x)\) is difficult to compute.
Choose a simpler candidate \(Q\) with probabilities \(q(z)\), supported
where \(p(x,z)>0\). Bayes' formula gives
\begin{align}
  D_{\mathrm{KL}}(Q\Vert P(\,\cdot\mid x))
  &=\sum_z q(z)\bigl[\log q(z)-\log p(x,z)\bigr]+\log p(x).
  \label{eq:variational-kl}
\end{align}
The evidence \(p(x)\) does not depend on the choice of \(Q\).
We can therefore search for the best approximation in the chosen family by
minimizing the sum in~\eqref{eq:variational-kl}, without knowing \(p(x)\).
This is the basic idea of \emph{variational inference}. It does not
guarantee that the chosen family contains the true posterior.

\subsection{Confidence is not correctness}

The entropy of a model's predicted distribution measures how concentrated
its own predictions are. For a prediction \(q=0.99\),
\[
  H(P_{0.99})\approx0.0560,
  \qquad -\log(1-0.99)\approx4.6052.
\]
The entropy is small, but if the observed label is zero, the log loss is
large. Predictive entropy alone does not establish accuracy or calibration.
This differs from the population cross-entropy, which averages prediction
losses using the distribution of the actual outcomes.

\subsection{A caution for continuous models}

For densities \(f\) and \(g\), the corresponding KL formula is
\[
  D_{\mathrm{KL}}(f\Vert g)
  =\int f(x)\log\frac{f(x)}{g(x)}\,dx,
\]
with the same support requirement and nonnegativity.
When the integrals are finite, the identity with cross-entropy remains
valid if entropy is replaced by \(-\int f\log f\), called
\emph{differential entropy}. This quantity can be negative and changes
under a change of measurement units. A density is not a point probability.

In particular, the empirical distribution of finitely many real-valued
observations places positive mass at individual sample points; a
continuous model places zero probability at each such point.
The KL from that atomic empirical distribution to the continuous model
is infinite. Identity~\eqref{eq:mle-empirical-ce} therefore applies to
finite labels. For continuous MLE, use the density likelihood from
Chapter~\ref{chap:estimation}.

\section{Score and Fisher information}
\label{sec:fisher-information}

Entropy describes uncertainty about an outcome. We now ask how precisely
the data determine an unknown parameter. Start with the likelihood peak
from Chapter~\ref{chap:estimation}: a broad peak leaves many nearby
parameter values with similar likelihood, while a narrow peak favors
a smaller range of values.

\subsection{Observed information at a fitted maximum}
\label{sec:observed-information}

Let \(\theta\) be a scalar parameter, and let \(f_\theta\) be the
observation PMF or density. For recorded independent observations
\(x_1,\ldots,x_n\), write
\[
  \ell_n(\theta)=\log L(\theta)
     =\sum_{i=1}^n\log f_\theta(x_i),\qquad
  S_n(\theta)=\ell_n'(\theta).
\]
The derivative \(S_n\) is the \emph{sample score}.
At a differentiable interior MLE \(\widehat\theta\),
\(S_n(\widehat\theta)=0\). A flat peak and a sharp peak both have zero
slope at their maximum. The second derivative distinguishes them.

\begin{definition}[Observed information]
For a twice differentiable log-likelihood, the observed information is
\begin{equation}
  J_n(\theta)=-\ell_n''(\theta).
  \label{eq:observed-information}
\end{equation}
It measures the downward curvature for the particular recorded dataset.
At an interior maximum \(J_n(\widehat\theta)\ge0\); away from a maximum,
observed information need not be nonnegative.
\end{definition}

Let \(\delta\) be a small displacement from the fitted parameter.
Since the linear term vanishes at the MLE, a quadratic approximation gives
\begin{result}[Likelihood drop near an interior MLE]
If the log-likelihood is well approximated by a quadratic near
\(\widehat\theta\), with \(J_n(\widehat\theta)>0\), then
\begin{align}
  \ell_n(\widehat\theta+\delta)-\ell_n(\widehat\theta)
     &\approx-\tfrac12 J_n(\widehat\theta)\delta^2,
       \label{eq:observed-log-drop}\\
  \frac{L(\widehat\theta+\delta)}{L(\widehat\theta)}
     &\approx\exp[-\tfrac12 J_n(\widehat\theta)\delta^2].
       \label{eq:observed-relative-likelihood}
\end{align}
\end{result}

For the same displacement \(\delta=0.1\), curvature \(100\) gives a
relative likelihood of approximately \(e^{-0.5}=0.607\), while
curvature \(4\) gives \(e^{-0.02}=0.980\). The first likelihood loses
much more of its height over the same parameter change.
At distance \(1/\sqrt{J_n(\widehat\theta)}\), the approximation drops
to \(e^{-1/2}\approx0.607\) of its maximum.
This is a local description of the fitted peak. Boundary MLEs need
separate treatment because their score need not vanish.

\subsection{Fisher information averages over possible samples}

Before collecting data, we do not know which likelihood peak we will
obtain. Keep a parameter value \(\theta\) fixed, repeatedly draw samples
from that model, and evaluate each sample's observed information at
the same \(\theta\). The average curvature is Fisher information.

\begin{definition}[Fisher information]
For a regular model satisfying the conditions below, the Fisher
information about \(\theta\) in a sample of size \(n\) is
\begin{equation}
  I_n(\theta)=\Expect_\theta[J_n(\theta)],
  \label{eq:fisher-expected}
\end{equation}
when this expectation exists. The expectation draws new samples under
\(\theta\); the parameter at which the curvature is evaluated stays fixed.
\end{definition}

Thus \(J_n\) describes the observed dataset, whereas \(I_n\) describes
the curvature expected under the model and sampling design.
For one observation this definition reads
\begin{equation}
  I_1(\theta)
     =-\Expect_\theta[\partial_\theta^2\log f_\theta(X)].
  \label{eq:fisher-curvature}
\end{equation}
The expectation in~\eqref{eq:fisher-expected} is not taken at a new
fitted maximum for each dataset. Refitting each time would be a
different calculation. We can evaluate the resulting model function
\(I_n(\theta)\) at \(\widehat\theta\) after obtaining an estimate.

\subsection{The equivalent squared-score formula}

To connect expected curvature to scores, assume an interior parameter
value, a common support not depending on \(\theta\), positive
\(f_\theta\) on that support, and differentiability allowing the
derivatives below to pass through the integral or sum. Assume also a
finite score second moment. These conditions exclude the moving
support of the uniform-endpoint model in Chapter~\ref{chap:estimation}.

\begin{definition}[Score of one observation]
The one-observation score is
\begin{equation}
  s_\theta(x)=\frac{\partial}{\partial\theta}\log f_\theta(x).
  \label{eq:score}
\end{equation}
A positive score means that a small increase in \(\theta\) increases
the log-likelihood of this observation.
\end{definition}

Normalization gives
\begin{equation}
  \Expect_\theta[s_\theta(X)]
  =\int\frac{\partial_\theta f_\theta(x)}{f_\theta(x)}
           f_\theta(x)\,dx
  =\frac{\partial}{\partial\theta}\int f_\theta(x)\,dx=0.
  \label{eq:score-mean}
\end{equation}
For a discrete model, use sums. Differentiating \(\log f_\theta\)
once more gives
\[
  \partial_\theta^2\log f_\theta
   =\frac{\partial_\theta^2 f_\theta}{f_\theta}
    -s_\theta^2.
\]
Multiply by \(f_\theta\) and integrate. The first term integrates to
the second derivative of \(1\), which is zero.

\begin{result}[Expected squared score and expected curvature]
Under the stated regularity conditions,
\begin{equation}
  I_1(\theta)=\Expect_\theta[s_\theta(X)^2]
            =\Var_\theta(s_\theta(X)).
  \label{eq:fisher}
\end{equation}
The variance equality uses the zero expected score in
\eqref{eq:score-mean}.
\end{result}

Squaring prevents positive and negative scores from cancelling.
For a small parameter change \(\delta\),
\(\log f_{\theta+\delta}(x)-\log f_\theta(x)
\approx\delta s_\theta(x)\). The squared-score expression measures
this sensitivity across possible observations.
The recorded sample score being zero at the MLE does not make Fisher
information zero: the expectation uses fresh samples at a fixed
parameter, whose scores vary.

For a Gaussian observation with known \(\sigma^2\),
\[
  s_\mu(x)=\frac{x-\mu}{\sigma^2},\qquad
  I_1(\mu)=\frac{\Expect_\mu[(X-\mu)^2]}{\sigma^4}
          =\frac1{\sigma^2}.
\]
A reading above \(\mu\) pushes toward a larger mean, and a reading below
it pushes toward a smaller mean. More observation noise gives less
information about the same mean.

\begin{result}[Independent observations add information]
For \(n\) i.i.d. observations,
\begin{equation}
  S_n(\theta)=\sum_{i=1}^n s_\theta(X_i),\qquad
  I_n(\theta)=\Expect_\theta[S_n(\theta)^2]=nI_1(\theta).
  \label{eq:fisher-additivity}
\end{equation}
\end{result}

The scores are independent and have mean zero, so their sum has variance
equal to the sum of their variances. Equivalently, the sample
log-likelihood is a sum, so its expected negative curvature adds.
With dependent observations, the score covariance terms need not vanish.

\subsection{When observed and expected information coincide}

For independent Gaussian measurements with known variance, completing
the square as in Section~\ref{sec:gaussian-estimation} gives
\[
  \widehat\mu=\bar x,\qquad
  \ell_n(\mu)=C-\frac{n}{2\sigma^2}(\mu-\bar x)^2,
\]
where \(C\) does not depend on the candidate mean \(\mu\).
The curvature is the same for every possible dataset:
\begin{equation}
  J_n(\mu)=-\ell_n''(\mu)=\frac{n}{\sigma^2}
         =I_n(\mu).
  \label{eq:gaussian-observed-expected}
\end{equation}
Consequently, the quadratic likelihood-drop formula is exact:
\[
  \frac{L(\bar x+\delta)}{L(\bar x)}
     =\exp\!\left[-\frac12 I_n(\mu)\delta^2\right].
\]
The equality of \(J_n\) and \(I_n\) here uses the known-variance
Gaussian mean model; it is not a general identity for every model.

\begin{figure}[H]
  \centering\InfoLikelihoodPlot
  \caption{Gaussian relative likelihoods with known \(\sigma=2\),
  centered at each sample's fitted mean. Here \(J_n=I_n=n/4\):
  the three curvatures are \(1.25,\ 5,\ 20\).
  The distance to relative height \(e^{-1/2}\) is \(1/\sqrt{I_n}\),
  so larger information means a narrower peak. These curves compare
  candidate means; their areas are not posterior probabilities.}
  \label{fig:information-likelihood}
\end{figure}

\subsection{Bernoulli information}

Now treat the population success probability \(p\) as unknown,
with \(0<p<1\). For \(Y\sim\operatorname{Bernoulli}(p)\),
\[
  \log f_p(y)=y\log p+(1-y)\log(1-p),
  \qquad
  s_p(y)=\frac y p-\frac{1-y}{1-p}
        =\frac{y-p}{p(1-p)}.
\]
At \(p=0.8\), a success gives score \(1.25\), while a failure gives
score \(-5\). Their weighted average is
\(0.8(1.25)+0.2(-5)=0\).
Using the Bernoulli variance already established in Chapter~2,
\begin{equation}
  I_1(p)=\frac{\Expect_p[(Y-p)^2]}{p^2(1-p)^2}
        =\frac1{p(1-p)},\qquad
  I_n(p)=\frac n{p(1-p)}.
  \label{eq:bernoulli-fisher}
\end{equation}

Observed information can differ from this expectation. Let
\(K=\sum_{i=1}^nY_i\) be the success count. Differentiating the sample
log-likelihood twice gives
\begin{equation}
  J_n(p)=\frac{K}{p^2}+\frac{n-K}{(1-p)^2}.
  \label{eq:bernoulli-observed}
\end{equation}
Unlike \(I_n(p)\), this quantity depends on which dataset was observed.
For example, evaluate two samples of size \(10\) at the same fixed
candidate probability \(p=0.8\):
\begin{center}
\begin{tabular}{@{}rrr@{}}
\toprule
Observed successes & \(J_{10}(0.8)\) & \(I_{10}(0.8)\)\\
\midrule
\(7\) & \(85.9375\) & \(62.5\)\\
\(9\) & \(39.0625\) & \(62.5\)\\
\bottomrule
\end{tabular}
\end{center}
The likelihood curvatures differ even though the model parameter used
for evaluation is the same. Averaging over all possible samples,
using \(\Expect_p K=np\), recovers Fisher information:
\[
  \Expect_p[J_n(p)]
  =\frac{np}{p^2}+\frac{n(1-p)}{(1-p)^2}
  =\frac{n}{p(1-p)}=I_n(p).
\]

\begin{figure}[H]
  \centering\InfoFisherPlot
  \caption{Information about the Bernoulli probability \(p\), per
  observation. It is smallest at \(p=1/2\). This is not the entropy
  curve: entropy measures uncertainty about the next label; Fisher
  information measures sensitivity to a change in the parameter.}
  \label{fig:bernoulli-fisher}
\end{figure}

There is no contradiction between maximum entropy and minimum Fisher
information at \(p=1/2\). A balanced label is most unpredictable, and the
sample proportion also has its largest absolute variance \(p(1-p)/n\)
there. Near zero, a small absolute error in \(p\) can still be a large
relative error. Moreover, Fisher information changes with the choice
of parameter and its units; it is not an intrinsic score for how useful
an observation is for every possible question.

\section{The Cramér--Rao lower bound}
\label{sec:cramer-rao}

Chapter~\ref{chap:estimation} compared the bias and variance of particular
estimators. Fisher information gives a lower bound for the variance of
\emph{every unbiased estimator} satisfying the regularity conditions.

\begin{theorem}[Scalar Cramér--Rao bound]
Let \(T=T(X_1,\ldots,X_n)\) estimate \(\theta\) from a regular i.i.d.
model as above, with \(0<I_n(\theta)<\infty\).
Suppose \(T\) has finite variance, \(\Expect_\theta T=\theta\)
throughout an open interval of parameter values, and differentiating
its expectation under the sum or integral is valid. Then
\begin{equation}
  \Var_\theta(T)\ge\frac1{I_n(\theta)}
                 =\frac1{nI_1(\theta)}.
  \label{eq:cramer-rao}
\end{equation}
\end{theorem}

\begin{proof}
Let \(g_\theta(x_1,\ldots,x_n)\) be the joint sample PMF or density.
Its log-derivative is the sample score \(S_n(\theta)\).
Differentiate the unbiasedness identity:
\[
  1=\frac{\partial}{\partial\theta}\Expect_\theta T
   =\int T(x)\,\partial_\theta g_\theta(x)\,dx
   =\Expect_\theta[T S_n(\theta)].
\]
The score has mean zero, so this last expectation equals
\(\Cov_\theta(T,S_n(\theta))\). The covariance bound from
Section~\ref{sec:covariance-correlation} gives
\[
  1=\Cov_\theta(T,S_n(\theta))^2
    \le\Var_\theta(T)\Var_\theta(S_n(\theta))
    =\Var_\theta(T)I_n(\theta).
\]
Dividing proves the claim.
\end{proof}

\begin{definition}[Efficient unbiased estimator]
An unbiased estimator is efficient at a parameter value if its variance
equals the Cramér--Rao bound there. Such equality need not be attainable
in every model or at every sample size.
\end{definition}

\subsection{Two estimators that attain the bound}

For \(X_i\sim\mathcal N(\mu,\sigma^2)\) independently, with known
\(\sigma^2\), the bound and the sample-mean variance are
\[
  \frac1{I_n(\mu)}=\frac{\sigma^2}{n}
                 =\Var_\mu(\bar X_n).
\]
The sample mean is unbiased and efficient. By contrast, the estimator
\(T=X_1\) is also unbiased but has variance \(\sigma^2\): it discards
the other \(n-1\) measurements.

\begin{figure}[H]
  \centering\InfoVariancePlot
  \caption{Known Gaussian variance \(\sigma^2=4\). The line \(4/n\)
  is both the Cramér--Rao bound and the exact variance of the sample
  mean. Points estimate that variance from \(30{,}000\) independent
  datasets at each \(n\), using seed \(20261020\). Using only \(X_1\)
  has variance \(4\), regardless of how many other observations exist.
  Both axes use logarithmic scales.}
  \label{fig:information-variance}
\end{figure}

For i.i.d. Bernoulli labels, the sample proportion
\(\widehat p=\bar Y_n\) is unbiased, and
\[
  \Var_p(\widehat p)=\frac{p(1-p)}n=\frac1{I_n(p)}.
\]
It too attains the bound for \(0<p<1\), at every sample size.
This exact variance equality does not make a Gaussian confidence
approximation exact; Section~\ref{sec:proportion-intervals} showed why
small counts and boundaries remain problematic.

\subsection{Why unbiasedness cannot be dropped}

In the Gaussian mean example, consider \(T=\bar X_n/2\).
It has smaller variance,
\[
  \Var_\mu(T)=\frac{\sigma^2}{4n}
             <\frac{\sigma^2}{n},
  \qquad
  \Expect_\mu T-\mu=-\frac\mu2.
\]
This does not violate the bound, because the estimator is biased.
Its mean squared error is
\[
  \Expect_\mu[(T-\mu)^2]=\frac{\sigma^2}{4n}+\frac{\mu^2}{4}.
\]
Whether this is better than the sample mean depends on \(\mu\).
This is the same shrinkage behavior encountered in the Gaussian MAP
example in Section~\ref{sec:map-estimation}. Even at \(\mu=0\), the
estimator is not unbiased throughout an open interval, as the theorem
requires.

\section{MLE uncertainty from information}
\label{sec:mle-information}

The Cramér--Rao bound is a statement about unbiased estimators.
An MLE need not be unbiased at a finite sample size. Nevertheless,
regular MLEs often approach the information bound as the sample grows.
The reason connects the score to the CLT from Chapter~\ref{chap:limit-theorems}.

Let \(\theta\) be the true parameter and \(\widehat\theta\) an interior
MLE, so \(\ell_n'(\widehat\theta)=0\).
Near \(\theta\), a linear approximation to the score gives
\[
  0\approx\ell_n'(\theta)
       +\ell_n''(\theta)(\widehat\theta-\theta).
\]
The score \(\ell_n'(\theta)=S_n(\theta)\) has mean zero and variance
\(nI_1(\theta)\). If the observed curvature is close to its expectation,
\(-\ell_n''(\theta)\approx nI_1(\theta)\), rearranging gives
\begin{equation}
  \widehat\theta-\theta
  \approx\frac{S_n(\theta)}{nI_1(\theta)}.
  \label{eq:mle-score-approximation}
\end{equation}
The CLT makes the numerator approximately Gaussian. Dividing by
\(nI_1(\theta)\) gives an error scale \(1/\sqrt{nI_1(\theta)}\).

\begin{result}[Asymptotic normality under explicit local conditions]
Assume the regular score has finite positive information, the MLE is
consistent and interior with probability tending to one, and
\[
  -\frac1n\ell_n''(\widetilde\theta_n)
       \xrightarrow{\mathbb P}I_1(\theta)
\]
for any sequence of points \(\widetilde\theta_n\) between \(\theta\) and
\(\widehat\theta\). Here the convergence notation is the one from
Chapter~\ref{chap:limit-theorems}. Then
\begin{equation}
  \sqrt{nI_1(\theta)}(\widehat\theta-\theta)
     \xrightarrow{d}\mathcal N(0,1).
  \label{eq:mle-normal-information}
\end{equation}
\end{result}

The mean value theorem gives the preceding score expansion with equality
when its second derivative is evaluated at such an intermediate point.
The sample-score CLT and the stated curvature convergence then give
\eqref{eq:mle-normal-information}. Consistency alone is not enough;
the local curvature condition is also needed.

The result justifies a Gaussian reference distribution with approximate
standard error
\[
  \widehat{\operatorname{SE}}(\widehat\theta)
  =\frac1{\sqrt{nI_1(\widehat\theta)}}.
\]
It does not, by itself, prove convergence of the actual variances:
that further claim needs control of the tails of the estimation errors.
When the approximation and substitution are valid, the interval formula
from Chapter~\ref{chap:confidence-intervals} becomes
\begin{equation}
  \widehat\theta\ \pm\
  z_{1-\alpha/2}\frac1{\sqrt{nI_1(\widehat\theta)}}.
  \label{eq:fisher-interval}
\end{equation}
The observed-information alternative uses
\(1/\sqrt{J_n(\widehat\theta)}\) when the fitted curvature is positive
and consistently estimates the same information.

For \(80\) successes in \(100\) independent Bernoulli observations,
\[
  \widehat p=0.8,\qquad
  nI_1(\widehat p)=\frac{100}{0.8(0.2)}=625,\qquad
  \widehat{\operatorname{SE}}=\frac1{\sqrt{625}}=0.04.
\]
The approximate \(95\%\) interval is
\[
  0.8\pm1.96(0.04)=[0.7216,\ 0.8784].
\]
This is the Wald interval already discussed in
Section~\ref{sec:proportion-intervals}. Fisher information explains its
standard error; it does not fix the boundary and small-sample
approximation problems.

\section{A local connection between KL and Fisher information}
\label{sec:local-kl-information}

KL compared two distributions; Fisher information described changing a
parameter. They meet when we compare nearby members of a regular family.
Fix \(\theta\), let \(\delta\) be a small parameter change, and define
\[
  F(\delta)=D_{\mathrm{KL}}(P_\theta\Vert P_{\theta+\delta})
  =\Expect_\theta[\log f_\theta(X)-\log f_{\theta+\delta}(X)].
\]
The expectation stays under \(P_\theta\) as \(\delta\) changes.
Under the differentiation assumptions used above,
\[
  F(0)=0,\qquad
  F'(0)=-\Expect_\theta[s_\theta(X)]=0,\qquad
  F''(0)=-\Expect_\theta[\partial_\theta^2\log f_\theta(X)]
        =I_1(\theta).
\]
A quadratic approximation with this value, slope, and curvature gives
\begin{result}[KL for a small parameter change]
For a sufficiently smooth regular family and small \(\delta\),
\begin{equation}
  D_{\mathrm{KL}}(P_\theta\Vert P_{\theta+\delta})
       \approx\frac12 I_1(\theta)\delta^2.
  \label{eq:kl-fisher-local}
\end{equation}
This is a local approximation, not a formula for arbitrary parameter
changes.
\end{result}

For Gaussians with the same known variance, the formula is exact:
\[
  \log\frac{f_\mu(x)}{f_{\mu+\delta}(x)}
    =\frac{-2\delta(x-\mu)+\delta^2}{2\sigma^2},
  \qquad
  D_{\mathrm{KL}}\!\left(
    \mathcal N(\mu,\sigma^2)\Vert\mathcal N(\mu+\delta,\sigma^2)\right)
    =\frac{\delta^2}{2\sigma^2}.
\]
The term linear in \(X-\mu\) disappears on taking the expectation.
Since \(I_1(\mu)=1/\sigma^2\), this equals
\(\tfrac12 I_1(\mu)\delta^2\) for every \(\delta\).

For Bernoulli probabilities near \(p=0.8\), the approximation is
\[
  D_{\mathrm{KL}}(P_{0.8}\Vert P_q)
  \approx\frac{(q-0.8)^2}{2(0.8)(0.2)}
  =3.125(q-0.8)^2.
\]
\begin{figure}[H]
  \centering\InfoLocalKLPlot
  \caption{Exact Bernoulli KL and its quadratic approximation around
  \(p=0.8\). They agree near \(q=0.8\), but the quadratic expression
  misses the asymmetric growth farther away.}
  \label{fig:local-kl-information}
\end{figure}

For independent samples, KL adds just as Fisher information does:
the log ratio of joint probabilities is the sum of the observation
log ratios, and taking the expectation gives \(n\) times the
one-observation KL. Thus nearby parameter changes have sample
divergence approximately \(nI_1(\theta)\delta^2/2\).
Setting \(|\delta|=1/\sqrt{nI_1(\theta)}\), the standard-error scale
from the previous section, makes that approximation equal to \(1/2\).
The same information quantity therefore controls both the separation
of nearby observation laws and the scale of regular estimation error.

\endgroup
